% --------------------------------------------------------------------
% --------------------------------------------------------------------
% Modelo de Trabalho Acadêmico utilizando classe repUERJ para elaboração
% de teses, dissertação e trabalhos monográficos em geral.
%
% Este arquivo está editado na codificação de caracteres UTF-8.
%
% As referencia estão baseadas no modelo bibtex e citação em autor-data
%
% Este modelo foi criado por Dr. Luís Fernando de Oliveira.
% Professor Adjunto do Departamento de Física Aplicada e Termodinâmica
% Instituto de Física Armando Dias Tavares
% Universidade do Estado do Rio de Janeiro - UERJ
%
% A classe repUERJ.cls foi criada a partir do código original 
% disponibilizado pelo grupo CódigoLivre (coordenado por
% Gerald Weber).
% Foram feitas adequações para implementação das normas de elaboração
% de teses e dissertações da UERJ.
%
% Os estilos repUERJformat.sty codificam os elementos pré-textuais e
% pós-textuais.
% O estilo repUERJpseudocode.sty codifica a elaboração de algoritmos
% utilizando um glossário desenvolvido por mim (Luís Fernando), o mesmo
% usado em meu curso de Física Computacional.
%
% Todo este material está disponível também no meu site
%      http://sites.google.com/site/deoliveiralf
%
% As normas da UERJ para elaboração de teses e dissertações pode ser
% obtidas no documento disponível no site
%      http://www.bdtd.uerj.br/roteiro_uerj_web.pdf
%
% Agradecimentos ao grupo da Rede Sirius - BDTD e à Biblioteca Setorial
% da Física.
% --------------------------------------------------------------------
% --------------------------------------------------------------------
%
\documentclass[a4paper,12pt,oneside,onecolumn,final,fleqn]{config/repUERJ}
% ---
% Pacotes fundamentais 
% ---
\usepackage[brazil,british]{babel}  % adequação para o português Brasil
\usepackage[utf8]{inputenc} % Determina a codificação utilizada
                            % (conversão automática dos acentos)
\usepackage{makeidx}        % Cria o índice
\usepackage{hyperref}       % Controla a formação do índice
\usepackage{indentfirst}    % Endenta o primeiro paragrafo de
                            % cada seção.
\usepackage{graphicx}       % Inclusão de gráficos
\graphicspath{images}
\usepackage{subfig}
\usepackage{amsmath}        % pacote matemático
\usepackage{amssymb}        % pacote extra
\usepackage{mathrsfs}       % pacote extra
\usepackage{tikz}	    % awesome drawings
\usetikzlibrary{calc}               % to get node width and height
\usetikzlibrary{positioning}        % to use the relative positioning commands
\usetikzlibrary{arrows.meta}
\usetikzlibrary{mindmap}
\usetikzlibrary{shapes.geometric}   % to use diamond shape
\usetikzlibrary{backgrounds}

% ---
% Pacote auxiliar para as normas da UERJ
% ---
\usepackage[frame=no,algline=yes,font=default]{config/repUERJformat}
\usepackage{config/repUERJpseudocode}
% ---
% Pacotes de citacoes
% ---
\usepackage[alf]{abntex2cite}

% ********************************************************************
% ********************************************************************
% Informações de autoria e institucionais
% ********************************************************************
% ********************************************************************

%---------------------------------------------------------------------
% Imagens pretextuais (precisam estar no mesmo diretório deste arquivo .tex)
%---------------------------------------------------------------------

\logo{images/logo_uerj_cinza.png}
\marcadagua{images/marcadagua_uerj_cinza.png}{1}{160}{255}

%---------------------------------------------------------------------
% Informações da instituição
%---------------------------------------------------------------------

\instituicao{Universidade do Estado do Rio de Janeiro}
            {Centro de Tecnologia e Ci\^{e}ncias} 
            {Instituto de Matem\'{a}tica e Estat\'{\i}stica} 
            {} 

%---------------------------------------------------------------------
% Informações da autoria do documento
%---------------------------------------------------------------------

\autor{Artur}
      {Chiaperini Grover}
      {A. C. G.}

      \titulo{Intelig\^{e}ncia Computacional Usando M\'{a}quinas de Boltzmann}
% se não for usar a quarta palavra chave, deixar o campo vazio: {}
\palavraschaves{m\'{a}quina de Boltzmann}
               {m\'{a}quina restrita de Boltzmann}
               {diverg\^{e}ncia de Kullback-Leibler}
               {entropia relativa}

\title{Computational Intelligence Using Boltzmann Machines}
\keywords{Boltzmann machine}
         {restricted Boltzmann machine}
         {Kullback-Leibler divergence}
         {relative entropy}

\orientador{Dra.}%
           {Roseli Suzi}{Wedemann}%
           {IME/CComp --- UERJ}

%coorientador é opcional
%\coorientador{Cargo Titulação} 
%             {Nome}{Sobrenome} 
%             {Unidade --- Instituição} 

%---------------------------------------------------------------------
% Grau pretendido (Doutor, Mestre, Bacharel, Licenciado) e Curso
%---------------------------------------------------------------------

\grau{Mestre}  
\curso{Curso}

% área de concentração é opcional
%\areadeconcentracao{área}

%---------------------------------------------------------------------
% Informações adicionais (local, data e paginas)
%---------------------------------------------------------------------

\local{Rio de Janeiro} 
\data{22}{Novembro}{2019} 

% ********************************************************************
% ********************************************************************
% Configurações de aparência do PDF final
% ********************************************************************
% ********************************************************************

% alterando o aspecto da cor azul
\definecolor{blue}{RGB}{41,5,195}
%\definecolor{apricot}{RGB}{251,206,177}

% informações do PDF
\hypersetup{%
  %backref=true,
  %pagebackref=true,
  %bookmarks=true,
  unicode=false,
  pdftitle={\UERJtitulo},
  pdfauthor={\UERJautor},
  pdfsubject={\UERJpreambulo},
  pdfkeywords={PALAVRAS}{CHAVES}{\chaveA}{\chaveB}{\chaveC}{\chaveD},
  pdfproducer={\packagename},       % producer of the document
  pdfcreator={\UERJautor},
  colorlinks=true,       % false: boxed links; true: colored links
  linkcolor=black,       % color of internal links blue
  citecolor=black,       % color of links to bibliography blue
  filecolor=black,       % color of file links magenta
  urlcolor=black,
  bookmarksdepth=4
}

% ********************************************************************
% ********************************************************************
% Início do documento
% ********************************************************************
% ********************************************************************
% ---
% compila o índice; se não for usar, comentar
% ---
\makeindex
% ---
% ********************************************************************
% ********************************************************************
\begin{document}
% ----------------------------------------------------------
% ELEMENTOS PRE-TEXTUAIS
% ----------------------------------------------------------
\frontmatter
% ----------------------------------------------------------
% Capa e a folha de rosto
% ----------------------------------------------------------
\capa%
\folhaderosto%
% ----------------------------------------------------------
% Inserir a ficha catalográfica
% ----------------------------------------------------------

% A biblioteca deverá providenciar a ficha catalográfica. Salve a ficha no formato PDF.
% Use o nome do arquivo PDF como argumento do comando. Exemplo: ficha catalográfica
% no arquivo 'ficha.pdf'

%\fichacatalografica{ficha.pdf}

% Enquanto não possuir a ficha catalográfica, use o comando sem argumentos...

\fichacatalografica{}

% ----------------------------------------------------------
% Folha de aprovação
% ----------------------------------------------------------
% membros da banca: máximo 6
\begin{folhadeaprovacao}
  \assinatura{Dra. Patr\'{\i}cia Nunes da Silva}{IME/CComp --- UERJ}
  \assinatura{Dra. Zochil Gonz\'{a}lez Arenas}{IME/CComp --- UERJ}
%  \assinatura{terceiro membro titular da banca}{instituição}
% suplente só é incluído se efetivamente substitui um titular
%  \assinatura{primeiro membro suplente da banca}{instituição}
%  \assinatura{segundo membro suplente da banca}{instituição}
%  \assinatura{terceiro membro suplente da banca}{instituição instituição instituição instituição }
\end{folhadeaprovacao}
% ----------------------------------------------------------
% Dedicatória
% ----------------------------------------------------------
%\pretextualchapter{Dedicatória}
%
%\vfill
%Texto da dedicatória
% ----------------------------------------------------------
% Agradecimentos
% ----------------------------------------------------------
%\pretextualchapter{Agradecimentos}
%
%Texto de agradecimento
% ----------------------------------------------------------
% Epigrafe (opcional)
% ----------------------------------------------------------
\pretextualchapter{}

\vfill\
\begin{flushright}
    If no one died, it is just another story to be told.
    
    [Daniel Mirolhaum]
\end{flushright}
% ----------------------------------------------------------
% RESUMO
% ----------------------------------------------------------
\pretextualchapter{Resumo}

\referencia%

Texto do resumo em português.\\

\imprimirchaves%
% ----------------------------------------------------------
% Abstract
% ----------------------------------------------------------
\pretextualchapter{Abstract}

\reference%

Abstract in English.\\

\printkeys%
% ----------------------------------------------------------
% Listas de ilustrações e tabelas
% ----------------------------------------------------------
\listadefiguras%
%\listadetabelas%
% ----------------------------------------------------------
% Outras listas
% ----------------------------------------------------------
%\listadealgoritmos%
% ----------------------------------------------------------
% Lista de abreviaturas e siglas
% ----------------------------------------------------------
%\pretextualchapter{Lista de abreviaturas e siglas}
\pretextualchapter{List of abbreviations and acronyms}

\abreviatura{BM}{Boltzmann Machine}
\abreviatura{RBM}{Restricted Boltzmann Machine}
\abreviatura{SNN}{Stochastic Neural Network}
\abreviatura{ANN}{Artificial Neural Network}
\abreviatura{KL}{Kullback-Leibler}

%\abreviatura{sigla1}{por extenso}
%\abreviatura{sigla2}{por extenso}
%\abreviatura{sigla3}{por extenso}
% ----------------------------------------------------------
% Lista de simbolos
% ----------------------------------------------------------
%\pretextualchapter{Lista de símbolos}

%\simbolo{simbolo1}{significado e/ou valor}
%\simbolo{simbolo2}{significado e/ou valor}
%\simbolo{simbolo3}{significado e/ou valor}
% ----------------------------------------------------------
% Sumario
% ----------------------------------------------------------
\sumario%
% ----------------------------------------------------------
% ELEMENTOS TEXTUAIS
% ----------------------------------------------------------
\mainmatter%
%=====================================================================
\chapter*{Introduction}%
%=====================================================================
\label{ch:intro}%
Nowadays, modelling intelligent complex systems uses two main paradigms, commonly referred to as Symbolic and Connectionist, as basic guidelines for achieving your goals of creating intelligent machines and understanding human cognition.
These two approaches depart from different positions, each advocating advantages over the other in reproducing intelligent activity.
The traditional symbolic approach argues that the algorithmic manipulation of symbolic systems is an appropriate context for modelling cognitive processes.
On the other hand, connectionists restrict themselves to brain-inspired architectures and argue that this approach has the potential to overcome the rigidity of symbolic systems by more accurately modeling cognitive tasks that can only be solved, in the best case, approximately. In addition, connectionist paradigm focus on parallel models which have simple and uniform processing elements interconnected~\cite{bib:sun2001}.

Years of experimentation with both paradigms lead us to the conclusion that the solution lies between these two extremes, and that the approaches must be integrated and unified.
In order to establish a proper link between them, much remains to be researched.
If in the 1980s the discussion of intelligence was placed at the distinct poles of the symbolists and connectionists, today the connectionists are divided by the reductionist arguments of the structuralists.
For this structuralist current, the failure of the symbolists was due to the fact that their models despised brain architecture, and therefore connectionism must continue to explore more deeply the structural aspects of the thinking organ.

The connectionist and structuralist aspects are approached, respectively, through the paradigm of artificial neural networks and realistic models of the brain, within the area called Computational Neuroscience.
Through our models, we investigate ancient questions of Artificial Intelligence regarding the understanding of computability aspects of the human mind.

In this project we will continue with the study and implementation of artificial neural network known as Boltzmann Machine and its derivatives models, such as the Restricted Boltzmann Machine, which have been used to solve artificial intelligence problems, in areas such as: computer vision, encoder problem, statistical pattern recognition, speech recognition, and combinatorial optimization~\cite{bib:hertz1991}.
Boltzmann machines is a branch of stochastic methods, where randomness is present.~\cite{bib:duda2000}, it is also used to solve \textit{searching} and \textit{learning} problems~\cite{bib:hinton2010}.

% automatic speech (or voice) recognition, image recognition and treatment, natural language processing, bioinformatics, among many others.

%Our previous experience, both in the development of research in the field of artificial neural networks and general distributed processing and its technological applications, as well as in the pursuit of realistic models of brain biology, allows us to mature in the same direction of multidisciplinary research.

%Texto da introdução. Texto, texto texto \cite{bib:Amado1991}, texto \citeonline{bib:Amado1991}. Texto \citeauthoronline{bib:Andrade1997} em \citeyear{bib:Andrade1997}, texto \citeauthor{bib:Andrade1997},  texto.

%=====================================================================
\chapter{Boltzmann Machines}%
%=====================================================================
\label{ch:bm}%
EXPLICAR OS ELEMENTOS DE PROBABILIDADE COM OS QUAIS ESTAMOS LIDANDO:%
QUEM S\~{A}O AS VARI\'{A}VEIS ALEAT\'{O}RIAS E COMO VAMOS DOMIN\'{A}-LAS NO TEXTO,
COMENTAR QUE ESTAMOS TRATANDO COM VARI\'{A}VEIS DISCRETAS, COMO VAMOS IDENTIFICAR AS PROBABILIDADES,$\ldots$

EXPLICAR O QUE OS \'{I}NDICES REPResentam, e as vari\'{a}veis.

EXPLICAR OS TERMOS $\omega$ e $\phi$. $\omega$ \'{E} para os pesos, e $\phi$ para o bias de cada unidade.

EXPLICAR MINHA NOTA\c{C}\~{A}O sobre o sobre-escrito entre par\^{e}nteses!!!


In this chapter will expose the Boltzmann Machine theory. 
A few considerations regarding the notation used in this exposition is required before stepping forward into the main content. 
A single random variables is denoted by $\mathrm{x}$.
A vector of random variables of size $n$ is represented by $\mathrm{\mathbf{x}} = (\mathrm{x}_{1}, \mathrm{x}_{2}, \dots, \mathrm{x}_{n})$, where each $\mathrm{x}_{i}$, $i \in \{1, 2, \dots, n\}$, represents a single unit of the network. 
The value a random variable can assume is represented by $x$. 
Assuming a discrete scenario, the probability of a random variable $\mathrm{x}_{i}$ of assuming a certain value $x_{i}$ is $P(\mathrm{x}_{i} = x_{i})$. 


\section{Boltzmann Machines}

Boltzmann Machines (BM) are a type of stochastic neural networks (SNN) where the connections between units, which are described by $\omega$, are symmetrical, i.e., $\omega_{ij} = \omega_{ji}$ [HERTZ].
This kind of stochastic neural networks are capable of learning internal representation and to model an input distribution.
Boltzmann Machines were named after the Boltzmann distribution.
Due to its stochatics behaviour, the probability of the state of the system to be found in a certain configuration is given by previous mentioned distribution [HERTZ].
According to [MONTUFAR, 2018], BM can be seen as an extension of Hopfield networks to include hidden units.

Boltzamann Machines have visible and hidden units.
The visible units are linked to the external world and they correspond to the components of an observation. On the other hand, the hidden units do not have any connection outside of the network and model the dependencies between the components of the observations [FISCHER, 2012].
In BM, there is no connection restriction, this means that every unit, visible or hidden, can be connected to every other unit as in a complete graph, this pattern is not mandatory as some of the connections may not exists depending on the network layout.
\begin{figure}[h!tbp]{\textwidth}
  \centering
  \caption{Boltzmann machine diagram.}%
  \label{fig:bm-diagram}%
  \tikzstyle{visible} = [draw, circle, fill=white]
  \tikzstyle{hidden} = [draw, circle, fill=gray!30]
  \begin{tikzpicture}
    \node[style=visible] (a) at (0, 2.5) {$\mathrm{x}_{1}$};
    % \node[style=visible] (b) at +(0: 1.5) {$\mathrm{x}_{2}$};
    % \node[style=hidden] (c) at + (60: 1.5) {$\mathrm{x}_{3}$};
    \node[style=visible] (b) at (-2.16, 1.25) {$\mathrm{x}_{2}$};
    \node[style=visible] (c) at (-2.16, -1.25) {$\mathrm{x}_{3}$};
    \node[style=visible] (d) at (0, -2.5) {$\mathrm{x}_{4}$};
    \node[style=hidden]  (e) at (2.16, -1.25) {$\mathrm{x}_{5}$};
    \node[style=hidden]  (f) at (2.16, 1.25) {$\mathrm{x}_{6}$};
    \foreach \from/\to in {a/b, a/c, a/d, a/e, a/f, b/c, b/d, b/e, b/f, c/d, c/e, c/f, d/e, d/f, e/f}
      \draw [thick] (\from) -- (\to);   
  \end{tikzpicture}
  \legend{Gray circles represent the hidden units of the Boltzmann Machine, while the white circles are the visible units.}
  \source{Author.}
\end{figure}


Training Boltzmann Machines means finding the right connection between the units.

Boltzmann Machines (BM) are stochastics neural networks with symmetric connections, i.e., $\omega_{ij} = \omega_{ji}$.
Boltzmann Machines use the Boltzmann distribution to determine the probability of the state of the system of the network.
BM ressambles the Hopfield networks with the inclusion of hidden units.
Finding the right connections between the hidden units without knowing it from the training patterns what the hidden units represent is part of the solving the Boltzmann Machine problem.

Units $x_{i}$ in BM are split into two kinds: visible and hidden units.
The visible units have connection to the outside world and are the units that receive the data input.
On the other hand, the hidden units do not have any connection to the outside of the network and they are resposible to find the data relation from the input.
In a BM, the connections between units can be complete or not.
Regardless of how the connections are, every connection in a BM is symmetric.

BM are made of stochastics units $x_{i}$.
Stochastics units are random variables that can assume a binary value with a certain probability. 
We will consider that a random variable $\mathrm{x}_{i}$ can assume a value $x_{i} \in \{0, 1\}$, i. e.,
\begin{equation}
  x_{i} =
    \begin{cases}
      1 \; \text{with probability} \; g(h_{i}) \\
      0 \; \text{with probability} \; 1 - g(h_{i})
    \end{cases},
  \label{eq:stochastic-unit-values}
\end{equation}
where the probability is given by
\begin{equation}
  g(h_{i}) = \frac{1}{1 + e^{-2 \beta h_{i}}},
  \label{eq:stochastic-unit-prob}
\end{equation}
and
\begin{equation}
  h_{i} = \sum_{j}\omega_{ij}x_{j}.
  \label{eq:stochastic-unit-input}
\end{equation}

Due to the symmetrical connections, there is an energy function give by
\begin{equation}
  H(\mathbf{x}) = - \sum_{i} \sum_{j} \omega_{ij}x_{i}x_{j} - \sum_{i} \phi_{i}x_{i},
  \label{eq:eq4}
\end{equation}
where $\mathbf{x} = (x_{1}, x_{2}, \dots, x_{n})$, and $n$ is equal to the number of units in the network, and the above energy function has minimum when there is a stable state characterised by
\begin{equation}
  x_{i} = sgn(h_{i}).
  \label{eq:eq5}
\end{equation}

The probability $P$ of finding the system in a given state $\mathbf{x}$ after the equilibrium is reached can be computed as follows:
\begin{equation}
  P(\mathbf{x}) = \frac{1}{Z} e^{-\beta H(\mathbf{x})},
  \label{eq:eq6}
\end{equation}
where
\begin{equation}
  Z = \sum_{\mathbf{x}'} e^{-\beta H(\mathbf{x}')}
  \label{eq:eq7}
\end{equation}
is the partition function.

The learning process of a Boltzmann Machine consists in ajusting the connections $\omega_{ij}$ in such a way that the state of the visible units have a particular desired probability distribution.

Let us identify the state of the visible units by an index $v$ and the state of the hidden units by an index $h$.
Considering a system which has N visible units and K hidden units, the whole system have $2^{N + K}$ possibilities of states in which it can be found.

The joint probability $P_{vh}$ is the probability of finding the visible and hidden units in the states $v$ and $h$, respectively.
This probabiblity measument is given by the Boltzmann distribution:
\begin{equation}
  P_{vh} = \frac{e^{-\beta H_{vh}}}{Z},
  \label{eq:eq8}
\end{equation}
where
\begin{equation}
  Z = \sum_{u} \sum_{k} e^{-\beta H_{uk}},
  \label{eq:eq9}
\end{equation}
and
\begin{equation}
  H_{vh} = - \sum_{i} \sum_{j} \omega_{ij} x^{(vh)}_{i} x^{(vh)}_{j} - \sum_{i} \phi_{i} x^{(vh)}_{i}.
  \label{eq:eq10}
\end{equation}

As metioned above, the problem a Boltzmann Machine is trying to solve is determining the connections $\omega_{ij}$ between units such that the visible units have a certain probability distribution.
In order to do that, we need to find the marginal probability of the state $v$ in which the visible units are found regardless of the state $h$ of the hidden units. The marginal probability $P_{v}$ is given by
\begin{equation}
  P_{v} = \sum_{h} P_{vh} = \sum_{h} \frac{e^{-\beta H_{vh}}}{Z}.
  \label{eq:marginal_prob}
\end{equation}

Although we know that $P_{v}$ is a function of the connections $\omega_{ij}$, and that this is the probability of finding the visible units in the state $v$. We want the states to have a certain probability $Q_{v}$, i.e., a desired probability.
This means that ideally we would like to match the empirical distribution of the data, even though we do not have access to the correct distribution, only to what the observed data has given us as an input to training the model.

One way to evaluate the difference between two probability distribution, for example, $P_{v}$ and $Q_{v}$, is using the Kullback-Leibler divergence, which can also be referred to relative entropy, $E$, which will be our cost function.
(EXPLICAR FUN\c{C}\~{A}O DE CUSTO e $D_{KL}$!!!).
\begin{equation}
  E = \sum_{v} Q_{v} \ln{\left(\frac{Q_{v}}{P_{v}}\right)}.
  \label{eq:DKL}
\end{equation}

The relative entropy $E$ has the property of always being equal or greater than zero.
It reaches zero only if $P_{v} = Q_{v}$, which means that we are able to retrieve the exactly probability distribution of the input data at the visible units.
\begin{equation}
  \begin{split}
    E & = \sum_{v} Q_{v} \ln{\left( \frac{Q_{v}}{P_{v}} \right)} \\
      & \geq \sum_{v} Q_{v} \left( 1 - \frac{P_{v}}{Q_{v}} \right) \\
      & = \sum_{v} \left( Q_{v} - P_{v} \right) \\
      & = \sum_{v} Q_{v} - \sum_{v} P_{v} = 1 - 1 \\
      & \Rightarrow E \geq 0.
  \end{split}
\end{equation}

From the gradient descent equation
\begin{equation}
  \Delta \omega_{ij} = -\eta \frac{\partial E}{\partial \omega_{ij}},
  \label{eq:gradient}
\end{equation}
where
\begin{equation}
  \begin{split}
    E & = \sum_{v} Q_{v} \ln{\left(\frac{Q_{v}}{P_{v}}\right)} \\
      & = \sum_{v} \left[ \ln{(Q_{v})} - \ln{(P_{v})} \right].
  \end{split}
  \label{eq:entropy1}
\end{equation}

In the following steps, we present the gradient descent derivation
\begin{equation}
  \begin{split}
    \Delta \omega_{ij} & = - \eta \frac{\partial E}{\partial \omega_{ij}} \\
                  & = - \eta \frac{\partial}{\partial \omega_{ij}} \left[ \sum_{v} Q_{v} \left( \ln{(Q_{v})} - \ln{(P_{v})} \right) \right] \\
                  & = \eta \frac{\partial}{\partial \omega_{ij}} \left[ \sum_{v} Q_{v} \ln{(P_{v})} \right] \\
                  & = \eta \sum_{v} Q_{v} \frac{\partial}{\partial \omega_{ij}} \left[ \ln{(P_{v})} \right] \\
                  \Rightarrow & \Delta \omega_{ij} = \eta \sum_{v} \frac{Q_{v}}{P_{v}} \frac{\partial P_{v}}{\partial \omega_{ij}}.
  \end{split}
  \label{eq:grad1}
\end{equation}

To continue with the computation of $\Delta \omega_{ij}$, we have to find the derivative of $\partial P_{v}/\partial \omega_{ij}$, from the marginal probability, equation~\ref{eq:marginal_prob},
\begin{equation}
  P_{v} = \frac{\sum_{h} e^{-\beta H_{vh}}}{\sum_{u} \sum_{k} e^{-\beta H_{u k}}},
  \label{eq:marginal_prob_expansion}
\end{equation}
thus the derivative of $P_{v}$ follows
\begin{equation}
  \begin{split}
    \frac{\partial P_{v}}{\partial \omega_{ij}} & = \frac{\partial}{\partial \omega_{ij}} \left[ \frac{\sum_{h} e^{-\beta H_{vh}}}{\sum_{u} \sum_{k} e^{-\beta H_{u k}}} \right] \\
    & = \frac{1}{\sum_{u} \sum_{k} e^{-\beta H_{u k}}} \sum_{h} (-\beta) e^{-\beta H_{vh}} \frac{\partial H_{vh}}{\partial \omega_{ij}} \\
    & - \sum_{h} e^{-\beta H_{vh}} \frac{1}{{\left( \sum_{u} \sum_{k} e^{-\beta H_{u k}} \right)}^{2}} \sum_{u} \sum_{k} e^{-\beta H_{u k}} (-\beta) \frac{\partial H_{u k}}{\partial \omega_{ij}}.
  \end{split}
  \label{eq:eq14}
\end{equation}

Following the equation~\ref{eq:eq14}, we need to compute the term $\partial H_{vh}/\partial \omega_{ij}$,
\begin{equation}
  \begin{split}
      \frac{\partial H_{vh}}{\partial \omega_{ij}} & = \frac{\partial}{\partial \omega_{ij}} \left[- \frac{1}{2} \sum_{m} \sum_{n} \omega_{mn} x^{(vh)}_{m} x^{(vh)}_{n} - \sum_{m} \phi_{mm} x^{(vh)}_{m} \right] \\
                                                   & = \frac{\partial}{\partial \omega_{ij}} \left[- \frac{1}{2} \sum_{m \neq i,j} \sum_{n \neq i,j} \omega_{mn} x^{(vh)}_{m} x^{(vh)}_{n} - \frac{1}{2} \omega_{ij} x^{(vh)}_{i} x^{(vh)}_{j} - \frac{1}{2} \omega_{ji} x^{(vh)}_{j} x^{(vh)}_{i} - \sum_{m} \phi_{m} x^{(vh)}_{m} \right],
  \end{split}
  \label{eq:eq15}
\end{equation}
as the connections between units are symmetric, i.e., $\omega_{ij} = \omega_{ji}$, then we can simplify equation~\ref{eq:eq15},
\begin{equation}
  \begin{split}
    \frac{\partial H_{vh}}{\partial \omega_{ij}} & = \frac{\partial}{\partial \omega_{ij}} \left[- \frac{1}{2} \sum_{m \neq i,j} \sum_{n \neq i,j} \omega_{mn} x^{(vh)}_{m} x^{(vh)}_{n} - \omega_{ij} x^{(vh)}_{i} x^{(vh)}_{j} - \sum_{m} \phi_{m} x^{(vh)}_{m} \right] \\
    & = \frac{\partial}{\partial \omega_{ij}} \left[- \frac{1}{2} \sum_{m \neq i,j} \sum_{n \neq i,j} \omega_{mn} x^{(vh)}_{m} x^{(vh)}_{n} \right] + \frac{\partial}{\partial \omega_{ij}} \left[- \omega_{ij} x^{(vh)}_{i} x^{(vh)}_{j} \right] + \frac{\partial}{\partial \omega_{ij}} \left[- \sum_{m} \phi_{m} x^{(vh)}_{m} \right] \\
    \Rightarrow & \frac{\partial H_{vh}}{\partial \omega_{ij}} = - x^{(vh)}_{i} x^{(vh)}_{j}.
  \end{split}
  \label{eq:eq16}
\end{equation}

Analagous to $\partial H_{vh}/\partial \omega_{ij}$, we have $H_{uk}$ derivative, which is
\begin{equation}
  \frac{\partial H_{uk}}{\partial \omega_{ij}} = - x^{(uk)}_{i} x^{(uk)}_{j}.
  \label{eq:eq17}
\end{equation}

Going back to equation~\ref{eq:eq14}, we can replace the derivatives of $H$, and solve the derivative of the marginal probability $P_{v}$,
\begin{equation}
  \begin{split}
    \frac{\partial P_{v}}{\partial \omega_{ij}} & = \frac{1}{Z} \sum_{h} e^{-\beta H_{vh}} (-\beta) (-x^{(vh)}_{i} x^{(vh)}_{j}) - \frac{1}{Z^{2}} \sum_{h} e^{-\beta H_{vh}} \sum_{u} \sum_{k} e^{-\beta H_{uk}} (-\beta) (-x^{(uk)}_{i} x^{(uk)}_{j}) \\
    & = \beta \left[ \sum_{h} \frac{e^{-\beta H_{vh}}}{Z} x^{(vh)}_{i} x^{(vh)}_{j} - \sum_{h} \frac{e^{-\beta H_{vh}}}{Z} \sum_{u} \sum_{k} \frac{e^{-\beta H_{uk}}}{Z} x^{(uk)}_{i} x^{(uk)}_{j}  \right] \\
    & = \beta \left[ \sum_{h} P_{vh} x^{(vh)}_{i} x^{(vh)}_{j} - P_{v} \sum_{u} \sum_{k} P_{uk} x^{(uk)}_{i} x^{(uk)}_{j} \right] \\
    \Rightarrow & \frac{\partial P_{v}}{\partial \omega_{ij}} = \beta \left[ \sum_{h} P_{vh} x^{(vh)}_{i} x^{(vh)}_{j} - P_{v} \langle x_{i} x_{j} \rangle  \right].
  \end{split}
  \label{eq:eq18}
\end{equation}

Given the derivative of $P_{v}$ in relation to $\omega_{ij}$, we can compute the learning term $\Delta \omega_{ij}$, from equation~\ref{eq:grad1},
\begin{equation}
  \begin{split}
    \Delta \omega_{ij} & = \eta \sum_{v} \frac{Q_{v}}{P_{v}} \beta \left[ \sum_{h} P_{vh} x^{(vh)}_{i} x^{(vh)}_{j} - P_{v} \langle x_{i} x_{j} \rangle \right]  \\
    & = \eta \beta \left[ \sum_{v} \sum_{h} \frac{Q_{v}}{P_{v}} P_{vh} x^{(vh)}_{i} x^{(vh)}_{j} - \sum_{v} \frac{Q_{v}}{P_{v}} P_{v} \langle x_{i} x_{j} \rangle \right] \\
    & = \eta \beta \left[ \sum_{v} \sum_{h} Q_{v} \frac{P_{vh}}{P_{v}} x^{(vh)}_{i} x^{(vh)}_{j} - \sum_{v} Q_{v} \langle x_{i} x_{j} \rangle \right] \\
    & = \eta \beta \left[ \sum_{v} \sum_{h} Q_{v} P_{h|v} x^{(vh)}_{i} x^{(vh)}_{j} - \langle x_{i} x_{j} \rangle \right].
  \end{split}
  \label{eq:learning_term2}
\end{equation}

In the above derivation, we have used the following relations,
\begin{equation}
  P_{h|v} = \frac{P_{vh}}{P_{v}},
  \label{eq:conditional_prob1}
\end{equation}
which is the conditional probability equation.
In our scenario this equation means that the probability distribution of the hidden units in state $h$ given the state $v$ of the visible units is the joint probability distribution of both states, $v$ and $h$, divided by the marginal probability distribution of the visible units in state $v$.

The second term in equation~\ref{eq:learning_term2}, is the average of units $i$ and $j$ over all combinations of states $v$ and $h$ of the system. 
In other words, we would have to compute all possible combination of states of visible and hidden units, $v$ and $h$, and then average over the specific units $i$ and $j$. 

The first term can be simplified by
\begin{equation}
  \sum_{v} Q_{v} \sum_{h} P_{h|v} x^{(vh)}_{i} x^{(vh)}_{j} = \sum_{v} Q_{v} \langle x_{i} x_{j} \rangle^{(v)} = \langle \langle x_{i} x_{j} \rangle^{(v)} \rangle.
  \label{eq:clamped_term}
\end{equation}

Then equation~\ref{eq:learning_term2} becomes
\begin{equation}
  \Delta \omega_{ij} = \eta \beta \left[ \langle \langle x_{i} x_{j} \rangle^{(v)} \rangle_{clamped} - \langle x_{i} x_{j} \rangle_{free} \right],
  \label{eq:learning_term3}
\end{equation}
it is important to notice that the subscripts \textit{clamped} means that we have to fix a certain $v$ state on the visible units otherwise the second term in the equation does not have a reference. 
On the other hand, the subscript \textit{free} identify the $\ldots$



%=====================================================================
\chapter{EXTRA STUFF}
%=====================================================================


\section{New Ideas}
\subsection{ANSARI}

The energy computed in Boltzmann Machines is a measure of self-consistency of the system.
Randomness is fundamental to avoid local optima.
Exponential schedule for annealing instead of linear approach, does not have the drawback of premature freezing.
Premature freezing means that the system reaches minimum temperature and becomes deterministic before a global minimum has been reached.

Boltzmann machines can be seem as a system where each unit represents a hypothesis and each connection a constrain among hypothesis.

The ability of the system to escape local optimum depends on the path chosen for the temperature, the annealing schedule. The success of simulated annealing depends on the proper choice of the annealing schedule.
There is a need to have an annealing schedule where the temperature never gets stuck at a particular value, but varies dynamically in such a way so to increase the system optimality.


\subsection{SMOLENSKY}

\begin{equation}
    \label{eq:smolensky-weights}
    W_{i \alpha} = {\left( k_{\alpha} \right)}_{i} \frac{\sigma_{\alpha}}{|\mathbf{k}_{\alpha}|},
\end{equation}
where $W_{i \alpha}$ is the weight between atom $\alpha$ and feature $i$, $\mathbf{k}_{\alpha}$ is the connection vector of the atom $\alpha$, ${\left( k_{\alpha} \right)}_{i}$ is the connection vector to a unit $i$, and $\sigma_{\alpha}$ is the strenght of the atom $\alpha$.

The harmony function assigns a real number $H_{\mathbf{K}}(\mathbf{r}, \mathbf{a}) $ to each state of the system, which is determined by a pair $(\mathbf{r}, \mathbf{a})$, where $\mathbf{r}$ is the representation vector and $\mathbf{a}$ the activation vector.
\begin{equation}
    \label{eq:smolensky-harmonyfunction}
    H_{\mathbf{K}}(\mathbf{r}, \mathbf{a}) = \sum_{\alpha} \sigma_{\alpha} a_{\alpha} h_{\kappa}(\mathbf{r}, \mathbf{k}_{\alpha}),
\end{equation}
where $a_{\alpha}$ is the connection to atom $\alpha$, and $h_{\kappa}(\mathbf{r}, \mathbf{k}_{\alpha})$ is the harmony contribution by activating the atom $\alpha$ given the current representation vector $\mathbf{r}$.
\begin{equation}
    \label{eq:smolensky-harmonyatom}
    h_{\kappa}(\mathbf{r}, \mathbf{k}_{\alpha}) = \frac{\mathbf{r} \cdot \mathbf{k}_{\alpha}}{|\mathbf{k}_{\alpha}|} - \kappa
    = \frac{\langle \mathbf{r}, \mathbf{k}_{\alpha} \rangle}{|\mathbf{k}_{\alpha}|} - \kappa.
\end{equation}

Putting equation~(\ref{eq:smolensky-harmonyatom}) into equation~(\ref{eq:smolensky-harmonyfunction}),
\begin{equation}
  \label{eq:smolensky-harmonyfunction2}
  \begin{split}
      H_{\mathbf{K}}(\mathbf{r}, \mathbf{a}) & =  \sum_{\alpha} \sigma_{\alpha} a_{\alpha} \left[ \frac{\langle \mathbf{r}, \mathbf{k}_{\alpha} \rangle}{|\mathbf{k}_{\alpha}|} - \kappa \right] \\
      & = \sum_{\alpha} a_{\alpha} \left[ \frac{\sigma_{\alpha} \langle \mathbf{r}, \mathbf{k}_{\alpha} \rangle}{|\mathbf{k}_{\alpha}|} - \sigma_{\alpha} \kappa \right].
  \end{split}
 \end{equation}
Considering the inner product
\begin{equation}
    \label{eq:smolensky-innerproduct}
    \langle \mathbf{r}, \mathbf{k}_{\alpha} \rangle = \sum_{i} r_{i} {\left( k_{\alpha} \right)}_{i},
\end{equation}
and replacing equation~(\ref{eq:smolensky-innerproduct}) into equation~(\ref{eq:smolensky-harmonyfunction2}),
\begin{equation}
  \label{eq:smolensky-harmonyfunction3}
  \begin{split}
      H_{\mathbf{K}}(\mathbf{r}, \mathbf{a}) & = \sum_{\alpha} a_{\alpha} \left[ \frac{\sigma_{\alpha}}{|\mathbf{k}_{\alpha}|}  \sum_{i} r_{i} {\left( k_{\alpha} \right)}_{i} - \sigma_{\alpha} \kappa \right] \\
  & = \sum_{\alpha} a_{\alpha} \left[ \sum_{i} r_{i} {\left( k_{\alpha} \right)}_{i} \frac{\sigma_{\alpha}}{|\mathbf{k}_{\alpha}|} - \sigma_{\alpha} \kappa \right].
  \end{split}
\end{equation}

Equation~(\ref{eq:smolensky-harmonyfunction3}) we can be reorder thus it becomes a function of the weights $W_{i \alpha}$, equation~(\ref{eq:smolensky-weights}),
\begin{equation}
    \label{eq:smolensky-harmonyfunctionfinal}
    H_{\mathbf{K}}(\mathbf{r}, \mathbf{a}) = \sum_{\alpha} \sum_{i} r_{i} W_{i \alpha} a_{\alpha} - \sum_{\alpha} a_{\alpha} \sigma_{\alpha} \kappa,
\end{equation}
where it is possible to see the harmony function as a function of the representation vector $\mathbf{r}$, the activation vector $\mathbf{a}$ just as in equation~(\ref{eq:smolensky-harmonyfunction}), but also as a function of the weights $W_{i \alpha}$ and the bias $\gamma_{\alpha} = \sigma_{\alpha} \kappa$ which is connected to each atom.
Both weights and bias carry the information of the atom strength $\sigma_{\alpha}$, which is related to the frequency of the atom $\alpha$: the more a particular atom $\alpha$ is activated, the more its frequency, thus increasing the importance of this atom to describe a feature $i$ from the representation vector $\mathbf{r}$.


The maximum-likelihood completion function $c(\iota)$ determined by the probability distribution $p$ is defined by
\begin{equation}
    \label{eq:smolensky-maxcompletion}
    \begin{split}
        c(\iota) = \{\mathbf{r} \in \mathbf{R} \: | \: \text{for some} \: \mathbf{a} \in \mathbf{A}, \text{and all} \: (\mathbf{r}', \mathbf{a}') \in \mathbf{R} \times \mathbf{A} \\
        \text{such that} \: \mathbf{r}' \supset \iota: p(\mathbf{r}, \mathbf{a}) \geqslant p(\mathbf{r}', \mathbf{a}') \},
   \end{split}
\end{equation}
where $\mathbf{r}$ is the completion of a point $\iota$.

In information theory, the concept of entropy (missing information) can be interpreted as the average of uncertainty of a variable outcome, similarly in physics, entropy is relate to the randomness of a system or the lack of information about it.
Entropy is a function of the probability distribution of a random variable $x$:
\begin{equation}
    \label{eq:smolensky-entropy}
    S(P(x)) = - \sum_{x \in X} P(x) \ln{(P(x))}.
\end{equation}

Theorem $1$: \textit{competence}.

A cognitive system should performe completion according to Boltzmann distribution for statistical extrapolation and inference.
Considering an experiment, each pattern can be observed with a particular frequency.
Assuming the probability distribution of the experiment $R$ is consistent with the Boltzmann distribution $P$ with harmony function given by equation~(\ref{eq:smolensky-harmonyfunctionfinal}), the maximum-likelihood completions of this distribution are the same as
\begin{equation}
    \label{eq:smolensky-teo1}
    P(\mathbf{r}, \mathbf{a}) \propto e^{H(\mathbf{r}, \mathbf{a})}.
\end{equation}

Theorem $2$: \textit{realizability}.

Nodes are updates stochastically, every node within the same layer is updated in parallel, and then every node in the other layer is updated in parallel.
During the update process the computational temperature is lowered via simulated annealing.


\subsection{AGGARWAL}

RBM\: the probabilistic states of the network are learned for a set of inputs.
This network models the joint probability distribution os the observed attributes together with some hidden attributes.
The connections are undirected because, RBM are designed to learn te probabilistic relationship rather than input-output mapping.
RBM are probabilistic models that create latent representations of the underlying data points.
It is similar to an autoencoder, but the latent representation is created stochastically instead of deterministically.

Boltzmann machines use probabilistic states to represent Bernoulli distributions of the binary attributes.
It contains both visible and hidden states.

Hopfield: given the weights and biases in the network at any particular time, it is possible to find a local energy minimum in terms of the states by repeatedly using the update rule:
\begin{equation}
    \label{eq:aggarwal-hopfield-update}
    s_{i} =
    \begin{cases}
         1 \; \text{if} \; \sum_{j : j \ne i} \omega_{ij} s_{j} + b_{i} \geq 0 \\
         0 \; \text{otherwise}
     \end{cases}.
\end{equation}

For a particular set of parameters $(\omega_{ij}, b_{i})$, the Boltzmann machine defines a probability distribution over various state configurations.
In Boltzmann machines, the dependence between all pairs of states is undirected: the visible states depend as much on the hidden states as the hidden states depend on visible states.
Boltzmann machine iteratively samples the states using a conditional distribution generated from the state values in the previous iteration unil thermal equilibrium is reached.

Learnings weights in the Boltzmann machine requires to maximize the likelihood of the specific training dataset observed.
It is also possible to maximize the log-likelihood, by taking the $\ln$ of the probability distribution.
\begin{equation}
    \label{eq:aggarwal-log-likelihood}
    \ln(P(s)) = - E(s) - \ln(Z)
\end{equation}

To maximize the log-likelihood, equation~(\ref{eq:aggarwal-hopfield-update}) is derived in terms of the weights $\omega_{ij}$, then it follows
\begin{equation}
    \label{eq:aggarwal-log-likelihood-max}
    \frac{\partial \ln(P(s))}{\partial \omega_{ij}} = {\langle s_{i}, s_{j} \rangle}_{data} - {\langle s_{i}, s_{j}
    \rangle}_{model}.
\end{equation}

Weights' update rule rely on two terms: ${\langle s_{i}, s_{j} \rangle}_{data}$ which represents the correlations between the states of the nodes $i$ and $j$ when the visible vectors are fixed to a vector in the training data, and the hidden states are allowed to vary; the second term, ${\langle s_{i}, s_{j} \rangle}_{model}$ is an average product of $s_{i}$ and $s_{j}$ from the model-centric samples obtained from Gibbs sampling. 
\begin{equation}
    \label{eq:aggarwal-bm-weight-update}
    \omega_{ij} \leftarrow \omega_{ij} \alpha ({\langle s_{i}, s_{j} \rangle}_{data} - {\langle s_{i}, s_{j} \rangle}_{model}),
\end{equation}

The subtraction of the model-centric correlations (second term) is required in order to remove the effect of the partition function within the expression of the log probability.
Monte Carlo sampling requires a large amount of samples to reach thermal equilibrium.
Restricted Boltzmann machine is a simplification of the Boltzmann machine.

The probability of each hidden unit taking on the value $1$ is:
\begin{equation}
    \label{eq:aggarwal-rbm-hidden-prob}
    P(h_{j} | \mathbf{v}) = \frac{1}{1 + e^{-b^{(h)}_{j} - \sum_{i} v_{i} \omega_{ij}}},
\end{equation}
where energy gap $\Delta E_{j}$ of a hidden unit $j$, between values $h_{j} = 0$ and $h_{j} = 1$ is
\begin{equation}
    \label{eq:aggarwal-rbm-energy-gap-hidden}
    \Delta E_{j} = b_{j} + \sum_{i} v_{i} \omega_{ij},
\end{equation}
analogously we have the probability of each visible unit.

Equation~(\ref{eq:aggarwal-rbm-hidden-prob}) can also be writen in terms of the sigmoid
\begin{equation}
    \label{eq:aggarwal-rbm-hidden-prob-sigmoid}
    P(h_{j} | \mathbf{v}) = \sigma \left( b_{j} + \sum_{i} v_{i} \omega_{ij} \right).
\end{equation}

For both the hidden and visible units, it is interesting to notice that the hidden unit $j$ depends only on the state of the visible units.
On the other hand, the visible unit $i$ depends only on the state of the hidden units.
This aspect makes it possible to update all hidden units in parallel and then all the visible units in parallel as well.

The weights encode the affinities between the visible and the hidden states.
Large positive weights implies the two states are likely to be on together.
Weights are initialized to a small value each.
Learning can be performed as follows:
\begin{itemize}
    \item \textit{Positive phase}: considering a mini-batch of training instances, for each element in this mini-batch, every hidden unit is updated according to equation~(\ref{eq:aggarwal-rbm-hidden-prob}).
    Once every hidden unit is updated, the average product between each pair of visible and hidden unit is computed, denoted as ${\langle v_{i}, h_{j} \rangle}_{pos}$.
    
    \item \textit{Negative phase}: for each training instance, it goes through a phase of Gibbs sampling after starting at with randomly initialized states.
    Is is achieved by repeatedly using equation~(\ref{eq:aggarwal-rbm-energy-gap-hidden}) and its version for the visible units.
    ${\langle v_{i}, h_{j} \rangle}_{neg}$ is computed after the values of $v_{i}$ and $h_{j}$ are at thermal equilibrium.
    
    \item Weights and biases can be updated as follows:
    \begin{equation}
        \label{eq:aggarwal-weights-bias-update}
        \begin{split}
            \omega_{ij} &\leftarrow \omega_{ij} + \alpha ({\langle v_{i}, h_{j} \rangle}_{pos} - {\langle v_{i}, h_{j} \rangle}_{neg}), \\
            b^{(v)}_{i} &\leftarrow b^{(v)}_{i} + \alpha ({\langle v_{i}, 1 \rangle}_{pos} - {\langle v_{i}, 1 \rangle}_{neg}), \\
            b^{(h)}_{j} &\leftarrow  b^{(h)}_{j} + \alpha ({\langle 1, h_{j} \rangle}_{pos} - {\langle 1, h_{j} \rangle}_{neg}), 
        \end{split}
    \end{equation}
    where $\alpha$ is the learning rate.
\end{itemize}

When visible unit $i$ and hidden unit $j$ are highly correlated, the weight $\omega_{ij}$ will be more positive, and the other way around, when visible unit $i$ and hidden unit $j$ are not correlated.
To obtain the negative samples, Monte Carlo sampling would be needed in order to reach thermal equilibrium.
It is possible to use an approximation of Monte Carlo sampling for only a short time and obtain the gradient, this method is called Contrastive Divergence.



The visible biases are initialized to $\ln(p_{i}/(1 - p_{i}))$, where $p_{i}$ is the fraction of data points in which the $i$th dimension takes on the value of $1$.
The hidden biases are initialized to $0$.



\subsubsection{Application of RBM}
RBM is only designed to learn probability distributions.
A mapping from input to output is required to deal with real-world problems.

The state of a node is a binary value sampled from the Bernoulli probabilities defined by equation~(\ref{eq:aggarwal-rbm-hidden-prob}).
On the other hand, the activation of a node in the associeted neural network is the probability value derived from the use of the sigmoid function in equation~(\ref{eq:aggarwal-rbm-hidden-prob-sigmoid}).

If RBM are used in a supervised application, a final phase of backpropagation is crucial.
In this case, the RBM are used as a pretraining for the supervised learning.




\subsection{Rubik's Code}

https://rubikscode.net/2018/10/01/introduction-to-restricted-boltzmann-machines/

Reason of statistical modelling and machine learning is to detect dependencies and connections between input variables.
In physics, energy is the capacity to do some sort of work.
The Energy-Based Models (EBMs) use the concept that the dependecies between variables can be determined by associating a scalar value to the whole energy of the system.
The associate value represents a measure of the probability that the system will be in a certain state.

\textbf{Inference:} process of making a prediction or a decision (when network is used).

\textbf{Learning:} process of finding an energy function that will associete low energies to correct values of the remaining variables and higher energy to incorrect values. (when parameters of the network are computed).

Boltzmann distribution states the probability of a state vector $\mathbf{s}$ is determined by the energy of that state vector relative to the energies of all possible binary state vectors.

The input do not need to be filled for every neuron; The Boltzmann machine will generate missing values.

Considering the neuron $i$ is given the chance to change its binary state, initially the total input of this neuron $i$ is computed by the weighted sum $z_{i}$ of the other active neurons connected to $i$ added by its own bias,
\begin{equation}
    \label{eq:rubik-input-weightedsum}
    z_{i} = b_{i} + \sum_{j} \omega_{ij} s_{j},
\end{equation}
where $b_{i}$ is the bias of neuron $i$, $s_{j}$ is the current state of the connected neurons to $i$, and $\omega_{ij}$ is the weight, i.e., the connection strength between neurons $j$ and $i$.
The probability that neuron $i$ will be active is given by:
\begin{equation}
    \label{eq:rubik-neuron-prob}
    p(s_{i} = 1) = \frac{1}{1 + e^{z_{i}}}.
\end{equation}

Neurons update theirselves randomly which means that the network will eventually reach a Boltzmann distribution, in which the probability of a state vector $\mathbf{v}$ is determined by the energy associated to that state vector relative to the energies of all possible binary state vectors:
\begin{equation}
    \label{eq:rubik-boltzmann-distri}
    P(\mathbf{v}) = \frac{e^{-E(\mathbf{v})}}{\sum_{\mathbf{u}} e^{-E(\mathbf{u})}}.
\end{equation}

The energy of the state of state vector $\mathbf{v}$ is defined as
\begin{equation}
    \label{eq:rubik-state-energy-bm}
    E(\mathbf{v}) = -\sum_{i} b_{i} s^{\mathbf{v}}_{i} - \frac{1}{2} \sum_{i} \sum_{j} s^{\mathbf{v}}_{i} \omega_{ij} s^{\mathbf{v}}_{j},
\end{equation}
where $s_{i}$ is the state assigned to neuron $i$ by state vector $\mathbf{v}$, analogously for neuron $j$.
The Boltzmann machine is resource-demanding as there can be connections between a neuron to every other neuron in the network.
Calculations can take a long time.

In the RBM architecture, there are no connections among visible units and no connections among hidden units, i.e., there are no visible-to-visible connection and no hidden-to-hidden connection.
There can be only connections between visible and hidden neurons, i.e., visible-to-hidden and the other way around as a connection between two units is symmetrical.

The energy function on RBM theory is modified to
\begin{equation}
    \label{eq:rubik-state-energy-rbm}
    E(\mathbf{v}, \mathbf{h}) = - \sum_{i} a_{i} v_{i} - \sum_{j} b_{j} h_{j} - \sum_{i} \sum_{j} v_{i} \omega_{ij} h_{j},
\end{equation}
where $a_{i}$ is the bias of the visible neuron $i$ in state $v_{i}$, $b_{j}$ is the bias of the hidden neuron $j$ in state $h_{j}$ and $\omeha_{ij}$ is the weight of the connection between visible neuron $i$ and hidden neuron $j$.

The joint probability of the intire system, when hidden neurons are in state vector $\mathbf{h}$ and visible neurons in state vector $\mathbf{v}$, is
\begin{equation}
    \label{eq:rubik-prob-rbm}
    P(\mathbf{v}, \mathbf{h}) = \frac{e^{-E(\mathbf{v}, \mathbf{h})}}{Z},
\end{equation}
where $Z$ is known as the partition function
\begin{equation}
    \label{eq:rubik-partition-func-rbm}
    Z = \sum_{\mathbf{u}} \sum_{\mathbf{k}} e^{-E(\mathbf{u}, \mathbf{k})},
\end{equation}
this partition function normalizes the probability and refers to the sum of all possible pairs of visible and hidden vectors of the system.

The RBM can be interpreted as a specialization architecture of the original Boltzmann machine, where a restriction to the neurons' connections have been put.

Given an input vector $\mathbf{v}$, the probability of a single hidden unit $j$ of being activated is 
\begin{equation}
    \label{eq:rubik-prob-hidden-active}
    P(h_{j} = 1 | \mathbf{v}) = \frac{1}{1 + e^{- \left( b_{j} + \sum_{i} v_{i} \omega_{ij} \right)}} = \sigma \left( b_{j} + \sum_{i} v_{i} \omega_{ij} \right),
\end{equation}
where $\sigma$ is the sigmoid function.
Analogous, the probability that a visible neuron $i$ is active given a hidden vector $\mathbf{h}$ is
\begin{equation}
    \label{eq:rubik-prob-visible-active}
    P(v_{i} = 1 | \mathbf{h}) = \frac{1}{1 + e^{- \left( a_{i} + \sum_{j} \omega_{ij} h_{j} \right)}} = \sigma \left( a_{i} + \sum_{j} \omega_{ij} h_{j} \right).
\end{equation}



\subsection{Contrastive Divergence Algorithm}

Here I try to compile the most information possible about Contrastive Divergence.
Each author is put in a separated subsubsection.
In the literature the contrastive divergence algorithm appears as $CD_{1}$, if it is a single iteration, or $CD_{k}$ for $k$ iterations.

\subsubsection{Aggarwal}

Contrastive divergence is a approximation approach to Monte Carlo sampling, it is usually applied as its fastest variant where only a single iteration is performed.

For a fixed training point of the dataset, i.e., a fixed visible state vector, for each hidden unit, its activation is computed, thus the hidden state vector is generated.
Then the visible units are generated again from the previous computed hidden state.
These values of the visible units are used as the sampled states stead of the ones obtained at thermal equilibrium. 
The hidden units are generated again from the last visible units computed. 

In the positive phase, only half an iteration of simply computing hidden states is used.
In the negative phase, there is at least one additional iteration to re-compute the hidden state.
This difference in the number of iterations is what causes the contrastive divergence between state distributions in among both phases.

An increase in the number of iterations causes the distribution to diverge from the data-conditioned states to what is proposed by the current weight vector. 
The value of $({\langle v_{i}, h_{j} \rangle}_{pos} - {\langle v_{i}, h_{j} \rangle}_{neg})$ quantifies the amount of contrastive divergence.
It is possible to increase the number of iterations $k$, which will improve the accuracy of the contrastive divergence.
Increasing $k$ will lead to better gradient at the cost of speed.
One approach is to progressively increase the value of $k$ as while applying the training.
\begin{itemize}
    \item weights as initialized to small values. In each iteration, only one additional step of contrastive divergence is used. One step is enough because the difference between weights are low in early iterations.
    \item as gradient descent improves near a better solution, increasing $k$ will improve accuracy. So two or three steps of CD can be applied. 
\end{itemize}

\subsubsection{Hinton}




%=====================================================================
%\chapter{T\'{i}tulo do cap\'{i}tulo 1}
%=====================================================================

%Texto do capítulo. Texto, texto, Figura~\ref{rotulo}. Texto Figura~\ref{outro.rotulo}\subref{subrotulo1}.

% \begin{figure ou table}[posição]{largura da figura}

%\begin{figure}[!ht]{6cm}
%  \caption{Título da figura.}\label{rotulo}
%  \includegraphics[width=\hsize]{images/logo_uerj_cor.jpg}
%  \legend{Texto da legenda.}
%  \source{Citação da fonte ou `O autor.'.}
%\end{figure}


%\begin{figure}[!ht]{11cm}
%  \caption{Título da figura.}\label{outro.rotulo}
%  \subfloat[][]{\label{subrotulo1}
%    \fbox{\includegraphics[width=0.45\hsize]{images/logo_uerj_cinza.png}}}\hfill
%  \subfloat[][]{\label{subrotulo2}
%    \fbox{\includegraphics[width=0.45\hsize]{images/marcadagua_uerj_cinza.png}}}\\
%  \subfloat[][]{\label{subrotulo3}
%    \fbox{\includegraphics[width=0.45\hsize]{images/logo_uerj_cor.jpg}}}\hfill
%  \legend{Texto da legenda. \subref{subrotulo1} Texto da imagem.
%          \subref{subrotulo2} Texto da imagem.
%          \subref{subrotulo3} Texto da imagem.}
%  \source{Citação da fonte ou `O autor'.}
%\end{figure}

%\begin{table}[!ht]{4cm}
%  \caption{Título da tabela.}\label{mais.rotulo}
%  \hfill\begin{tabular}{l|l}
%    \hline
%      X & Y\\
%    \hline
%      1,20 & 15,7\\
%      1,23 & 15,6\\
%      1,19 & 15,3\\
%      1,26 & 15,1\\
%      1,22 & 15,5\\
%      1,16 & 15,3\\
%      1,37 & 15,7\\
%    \hline
%  \end{tabular}\hfill
%  \legend{Texto da legenda.}
%  \source{Citação da fonte ou `O autor.'.}
%\end{table}

%=====================================================================
%\chapter{Título do capítulo 2}
%=====================================================================

%Texto do capítulo. Texto, texto Algoritmo~\ref{alg:modelo}. Texto.
%
%\begin{algorithm}[!ht]
%    \caption{Título do algoritmo.}\label{alg:modelo}
%    \begin{pseudocode}
%      \LinhaEmBranco%
%      \Documentacao%
%        \Titulo{Nome do algoritmo\\}
%        \Proposito{Propósito do algoritmo.\\}
%        \Metodo{Método utilizado no algoritmo.\\}
%        \Entradas{%
%          a, m: multiplicador e módulo\\
%          n0: semente\\
%          i: contador auxiliar\\
%        }
%        \Saidas{%
%          n: número aleatório\\
%        }
%        \Observacoes{Observações, restrições e requisitos.\\}
%        \Algoritmo{Identificação}
%          \Declarar{$a, m, i$}{numéricos}{}{}
%          \Declarar{$n0, n$}{numéricos}{}{}
%          \Ins{$m \leftarrow 13$}
%          \Ins{$n0 \leftarrow 1$}
%          \ParaDeAtePasso[para cada possível valor de `a']{$a$}{2}{$m-1$}{}
%            \Escrever{``a = '', $a$, ``: n = \{''}
%            \Ins[reinicia a geração com a semente n0]{$n \leftarrow n0$}
%            \ParaDeAtePasso{$i$}{0}{$m-1$}{}
%              \Ins[gerador de números aleatórios]{$n \leftarrow resto(a*n, m)$}
%              \SeEntao[se fim da sequencia $\dots$]{$n == n0$}
%                \Escrever{$n$,``\}''}
%                \Parar%
%              \Senao%
%                \Escrever{$n$}
%              \FimSe%
%            \FimPara%
%          \FimPara%
%        \Continua%
%    \end{pseudocode}
%\end{algorithm}
%
%\alglinenumbersoff%
%\begin{algorithm*}[!ht]
%    \caption{Título do algoritmo. (continuação)}
%    \begin{pseudocode*}
%        \LinhaEmBranco%
%        \Continuacao%
%%          \LinhaEmBranco
%          \Ins{$a \leftarrow 1$}
%          \Enquanto[comentário]{$a<10$}
%            \Escrever{$a$}
%            \Ins{$a \leftarrow a+1$}
%          \FimEnquanto%
%          \Ins{$a \leftarrow 1$}
%          \Repetir[comentário]
%            \Escrever{$a$}
%            \Ins{$a \leftarrow a+1$}
%          \AteQue{$a\ge10$}
%          \LinhaEmBranco%
%          \Ins{$a \leftarrow 1$}
%          \Fazer[comentário]
%            \Escrever{$a$}
%            \Ins{$a \leftarrow a+1$}
%          \Enquanto{$a<10$}
%        \FimAlgoritmo%
%      \FimDocumentacao%
%    \end{pseudocode*}
%\end{algorithm*}
%
%\vfill~\\

%=====================================================================
\chapter*{Future Development}%
%=====================================================================
\label{ch:future}%
\section*{To be studied$\ldots$}%
\label{ch:future:to-be-studied}%

Following up the current status of the project, the next steps includes:
\begin{enumerate}
  \item Finalization details of Boltzmann Machines, such as the update equations of the state of a unit and the description of the BM algorithm.
  \item Studying Restricted Boltzmann Machines and how its learning algorithm works.
  \item Dummy test to compare results of Hopfield, Boltzmann Machine and Restricted Boltzmann Machine
  \item Identify an interesting problem where the application of RBM is suitable.
\end{enumerate}

\subsection*{Chronogram}%
\label{ch:future:chronogram}

Aiming to conclude this masters dissertation the chronogram below seems to be the most coherent:
\begin{itemize}
  \item Finalization of both Boltzmann and Restricted Boltzmann Machines by the end of \textbf{January 2020}.
  \item Definition of the case to apply RBM by the end of \textbf{January 2020}.
  \item Implementation and testing of RBM algorithm to solve proposed task by the end of \textbf{May 2020}.
  \item Meanwhile keep with the development of the written text in parallel which will take up to about \textbf{June 2020}.
\end{itemize}


%=====================================================================
%\chapter*{Conclusion}
%=====================================================================

%Texto da conclusão.

% ----------------------------------------------------------
% ELEMENTOS POS-TEXTUAIS
% ----------------------------------------------------------
\backmatter%
%=====================================================================
% Referencias via BibTeX
%=====================================================================
\citeoption{abnt-options4.bib}
\bibliography{bibs/bibliografia.bib}
%=====================================================================

%=====================================================================
%\postextualchapter*{Glossário}
%=====================================================================
%\definicao{termo}{significado}
%\definicao{termo}{significado}
%\definicao{termo}{significado}
% ----------------------------------------------------------
% Apêndices (opcionais)
% ----------------------------------------------------------
% ---
% Inicia os apêndices
% ---
\appendix
%=====================================================================
\postextualchapter{Kullback-Leibler Divergence or Relative Entropy}%
%=====================================================================
\label{app:dkl}%
\section{Kullback-Leibler Divergence}%
\label{app:dkl:dkl}%

The Kullback-Leibler divergence or relative entropy is the measure of the difference between two probability distributions. 
The entropy computes how uncertainty events are. 
If a system is quite stable, the entropy will be close to zero and will not vary a lot, while if the system is quite unstable, then the entropy will be very large. 




%=====================================================================
\postextualchapter{Smolensky Word Completion Example}%
%=====================================================================
\label{app:make}%
\section{Smolensky MAKE Example}%
\label{app:smolensky:make}%

(CITAR SMOLENSKY) presents a word completion example with a Restricted Boltzmann Machine. 
In other words, the RBM has to complete missing or unknown letters of a given word. 
Based on previous example that had been shown to the RBM, the machine would complete the unknown information of the word with the most probable letter taking into account the other letters known in the word.

This example uses the word `MAKE'. (CITAR SMOLENSKY) calls the units on the hidden layer as `knowledge atoms', and the units on the visible layer as `feature processors'. In this appendix we will keep the nomenclature used by Smolensky.

The input to a completion task is provided by fixing some of the `feature processors' --- the known information --- while the unknown information is allowed to be updated. Fixed `features' are the units that are assigned the values $-1$ or $+1$, while the unknown `features' have value $0$. 





The entropy computes how uncertainty events are.
If a system is quite stable, the entropy will be close to zero and will not vary a lot, while if the system is quite unstable, then the entropy will be very large.
The Kullback-Leibler divergence, or relative entropy, is the measure of how different two separate probability distributions, $R$ and $P$, are from each other~\cite{bib:goodfellow2016}. Considering these probability distribution are both discrete and are defined on the same probability space $\chi$, the Kullback-Leibler divergence $D_{KL}(R || P)$ is given by
\begin{equation}
  \label{eq:app-dkl}
  D_{KL}(R || P) = \sum_{x \in \chi} R(x) \ln \left(\frac{R(x)}{P(x)} \right),
\end{equation}
the equation above can be read as the divergence from $P$ to $R$.

It is important to notice that this divergence measurement is not symmetric,
\begin{equation}
  \label{eq:app-dkl-nonsymmetric}
  \begin{split}
    D_{KL}(R||P) & \neq D_{KL}(P||R), \\
    \sum_{x \in \chi} R(x) \ln \left(\frac{R(x)}{P(x)} \right) & \neq \sum_{x \in \chi} P(x) \ln \left(\frac{P(x)}{R(x)} \right),
  \end{split}
\end{equation}
and that KL divergence, relative entropy, is always a non-negative value, 
\begin{equation}
  \label{eq:app-dkl-nonnegative}
  \begin{split}
    E & = D_{KL}(R||P) \\
    D_{KL}(R||P) & = \sum_{x \in \chi} R(x) \ln{\left( \frac{R(x)}{P(x)} \right)} \\
    & \geq \sum_{x \in \chi} R(x) \left( 1 - \frac{P(x)}{R(x)} \right) \\
    & = \sum_{x \in \chi} \left( R(x) - P(x) \right) \\
    & = \sum_{x \in \chi} R(x) - \sum_{x \in \chi} P(x) = 1 - 1 \\
    & \Rightarrow D_{KL}(R||P) \geq 0.
  \end{split}
\end{equation}

Notice that $D_{KL}(R||P)$ will only reach zero if and only if $P(x) = R(x)$, which means that we are able to retrieve the exactly desired probability distribution at the visible units from the input data.


\begin{figure}[h!tbp]{\textwidth}
    \centering
    \caption{EXAMPLE FIGURE FROM TIKZ}%
    \label{fig:smolensky-make}%
    \begin{tikzpicture}
        [node distance=1.3cm, on grid, >=stealth', bend angle=45, auto,
        every place/.style={minimum size=6mm, thick, draw=blue!75, fill=blue!20,}
        every transition/.style={thick, draw=black!75, fill=black!20},
        red place/.style={place, draw=red!75, fill=red!20},
        every label/.style={red}]

        \node [place, tokens=1] (w1) {};
        \node [place] (c1) [below=of w1] {};
        \node [place] (s) [below=of c1, label=above:$s\le 3$] {};
        \node [place] (c2) [below=of s] {};
        \node [place, tokens=1] (w2) [below=of c2] {};

        \node [transition] (e1) [left=of c1] {}
            edge [pre, bend left] (w1)
            edge [post, bend right] (s)
            edge [post] (c1);
        \node [transition] (e2) [left=of c2] {}
            edge [pre, bend right] (w2)
            edge [post, bend left] (s)
            edge [post] (c2);
        \node [transition] (l1) [right=of c1] {}
            edge [pre] (c1)
            edge [pre, bend left] (s)
            edge [post, bend right] node[swap] {2} (w1);
        \node [transition] (l2) [right=of c2] {}
            edge [pre] (c2)
            edge [pre, bend right] (s)
            edge [post, bend left] node {2} (w2);

        \begin{scope}[xshift=6cm]
            \node [place, tokens=1] (w1') {};
            \node [place] (c1') [below=of w1'] {};
            \node [red place] (s1') [below=of c1',xshift=-5mm] [label=left:$s$] {};
            \node [red place] (s2') [below=of c1',xshift=5mm] [label=right:$\bar s$] {};
            \node [place] (c2') [below=of s1',xshift=5mm] {};
            \node [place, tokens=1] (w2') [below=of c2'] {};

            \node [transition] (e1') [left=of c1'] {}
                edge [pre, bend left] (w1')
                edge [post] (s1')
                edge [pre] (s2')
                edge [post] (c1');
            \node [transition] (e2') [left=of c2'] {}
                edge [pre, bend right] (w2')
                edge [post] (s1')
                edge [pre] (s2')
                edge [post] (c2');
            \node [transition] (l1') [right=of c1'] {}
                edge [pre] (c1')
                edge [pre] (s1')
                edge [post] (s2')
                edge [post, bend right] node[swap] {2} (w1');
            \node [transition] (l2') [right=of c2'] {}
                edge [pre] (c2')
                edge [pre] (s1')
                edge [post] (s2')
                edge [post, bend left] node {2} (w2');
        \end{scope}

        \begin{scope}[on background layer]
            \node (r1) [fill=black!10,rounded corners,fit=(w1)(w2)(e1)(e2)(l1)(l2)] {};
            \node (r2) [fill=black!10,rounded corners,fit=(w1')(w2')(e1')(e2')(l1')(l2')] {};
        \end{scope}

        \draw [shorten >=1mm, -to, thick, decorate,
               decoration={snake, amplitude=.4mm, segment length=2mm,
                           pre=moveto, pre length=1mm, post length=2mm}]
              (r1) -- (r2) node [above=1mm, midway, text width=3cm, align=center]
              {replacement of the \textcolor{red}{capacity} by \textcolor{red}{two places}};
        
    \end{tikzpicture}   
    \legend{No legend}
    \source{PGFMANUAL.pdf}
\end{figure}









%\section{Primeira seção}
%
%Texto da primeira seção.
%
%\subsection{Primeira subseção}
%
%Texto da primeira subseção.
%
%\subsubsection{Primeira subsubseção}
%
%Texto da primeira subsubseção.
%=====================================================================
%\postextualchapter{Segundo apêndice}
%=====================================================================
%\section{Primeira seção}
%
%Texto da primeira seção.
%
%\subsection{Primeira subseção}
%
%Texto da primeira subseção.
%
%\subsubsection{Primeira subsubseção}
%
%Texto da primeira subsubseção.
% ----------------------------------------------------------
% Anexos (opcionais)
% ----------------------------------------------------------
% ---
% Inicia os anexos
% ---
\annex%
%=====================================================================
%\postextualchapter{Kullback-Leibler Divergence or Relative Entropy}
%=====================================================================
%\section{Primeira seção}
%
%Texto da primeira seção.
%
%\subsection{Primeira subseção}
%
%Texto da primeira subseção.
%
%\subsubsection{Primeira subsubseção}
%
%Texto da primeira subsubseção.
%=====================================================================
%\postextualchapter{Segundo anexo}
%=====================================================================
%\section{Primeira seção}
%
%Texto da primeira seção.
%
%\subsection{Primeira subseção}
%
%Texto da primeira subseção.
%
%\subsubsection{Primeira subsubseção}
%
%Texto da primeira subsubseção.
%---------------------------------------------------------------------
% INDICE REMISSIVO (relativo ao makeindex)
%---------------------------------------------------------------------
\printindex
%=====================================================================
\end{document}
